\documentclass[11pt]{article}
\usepackage[margin=1.15in]{geometry}
\usepackage{amsmath,amssymb,graphicx,xcolor,needspace}
\usepackage[hidelinks]{hyperref}
\newcommand{\Cov}{\mathsf C}
\newcommand{\T}{^{\mathsf T}}
\newcommand{\us}{\_\allowbreak}
\emergencystretch=3em
\begin{document}
\begin{center}
{\LARGE Intrinsic alignments in a joint shear forecast\par}
\vspace{.4em}
{\large From opposite-sign predictions to marginalized errors\par}
\vspace{.6em}
{\large November 2026\par}
\end{center}

\section{The scientific question and its executable calculation}
\paragraph{Why galaxy shapes need two contributions.}
A distant galaxy can appear elliptical because of its own shape and because
matter along the line of sight stretches its image. The gravitational part
is the shear. The intrinsic part need not be random: galaxies formed in the
same tidal environment can have correlated orientations. A change in this
alignment can imitate a change in the strength of matter fluctuations.
Our question is whether adding three-point information distinguishes those
changes after we account for their correlated measurement errors.

\paragraph{What the two papers establish.}
Gomes and collaborators' methods paper~\cite{methods} studies joint shear
inference with simulated DES Y3 observations, including covariance and
compression. Its reported 83\% improvement in the matter-density--amplitude
figure of merit belongs to that analysis. The TATT paper~\cite{tatt} studies
intrinsic-alignment contributions to three-point shear and aperture
skewness, using DES Y3 parameter choices. It finds examples where quadratic
alignment affects two- and three-point signals differently. It is a modeling
paper, not a DES Y6 posterior analysis. Our local Fisher calculations ask
related questions; reproducing their published posterior requires matching
their estimator, cuts, parameter space, priors and likelihood.

\paragraph{Where the student runs the calculation.}
The companion
\href{EXAMPLE_FISHER_IA.ipynb}{\texttt{EXAMPLE\us FISHER\us IA.ipynb}}
lives beside this note. It reads the project's numbered evaluate YAMLs,
builds their ordinary Cobaya models, and uses the shared CosmoLike Fisher
functions. The notebook contains the calculation and plotting cells;
there is no separate forecast driver. Open it from this directory or the
DES project directory in an activated Cocoa environment. Its first cell
finds the project and changes to the Cocoa directory so the YAML paths
retain their usual meaning. Read the markdown before running the costly
prediction and derivative cells. The committed notebook is executed on the
current inputs, and every forecast number quoted in this note is printed
by one of its cells.

\section{The observables and the alignment parameters}
\paragraph{A data vector is a list of measured numbers.}
The two-point entries are shear correlations, denoted by $\xi_+$ and $\xi_-$,
for different angular separations and pairs of source-redshift bins. The
three-point entries are third moments of the aperture mass. An aperture
mass first averages a shear field with a compensated angular filter; its
third moment then measures a non-Gaussian property of that filtered field.
The aperture radius is not the same quantity as a cut on triangle sides.
The present numbered examples supply 400 two-point entries and 80 aperture
entries, and they retain the supplied cross-covariance. The likelihood now
integrates its aperture means over the measured finite side bins: twenty
logarithmic bins from 0.5 to 80 arcminutes, cut to the nine whose lower
edge is at least 8 arcminutes, with first retained edge 8.151 arcminutes.
What remains open is narrower: the producer's exact bin-survival rule, the
63 opening-angle bin endpoints, and a covariance computed for the same
finite bins (see \nolinkurl{tests/studies/estimator_conventions.md}).

\paragraph{What NLA and TATT vary.}
The nonlinear alignment model, NLA, gives the linear tidal response an
amplitude $A_1$ and a redshift exponent $\eta_1$. TATT adds a quadratic tidal
response with amplitude $A_2$, its exponent $\eta_2$, and a density-weighting
coefficient $b_{\rm TA}$. A density-weighted contribution allows the
alignment to depend on where galaxies sample the matter field. The native
CosmoLike amplitude normalization and pivot are used throughout. Thus the
same parameter means the same physical input in two- and three-point
predictions. A change of convention belongs in a comparison to another
code, not in the middle of a forecast.

\paragraph{The redshift law.}
For illustration, the amplitude parameter's redshift factor is
\begin{equation}
 A_1(z)=A_1\left(\frac{1+z}{1+z_0}\right)^{\eta_1}.
\end{equation}
Here $z_0=0.62$ is the pivot. At the pivot the ratio inside parentheses is
one, so the factor equals $A_1$. The complete physical alignment coefficient
also contains the native matter-density, growth and normalization factors.
Changing $A_1$ reverses the linear contribution's sign; it does not reverse
every term of the complete TATT prediction. Products with even numbers of
alignment factors can retain their sign.

\paragraph{First make a finite comparison.}
The notebook evaluates a no-alignment mean, two NLA means, and two TATT
means with opposing amplitude choices. For TATT these choices are
\begin{equation}
 (A_1,A_2)=(-0.2,0.4),\qquad (A_1,A_2)=(0.2,-0.4).
\end{equation}
Cosmology, source distributions, source shifts and numerical settings stay
fixed. In the NLA cases both $A_2$ and $b_{\rm TA}$ vanish. This experiment
asks how two specified predictions differ; it does not optimize any other
parameter. Later derivatives ask a different, local question around each
numbered example's fiducial point.

\section{Why the cross-covariance changes the answer}
\paragraph{Write the matrix before using it.}
Partition the joint vector into its two- and three-point pieces:
\begin{equation}
 d=\begin{pmatrix}d_2\\d_3\end{pmatrix}.
\end{equation}
Its covariance has the corresponding blocks
\begin{equation}
 \Cov=\begin{pmatrix}C_{22}&C_{23}\\C_{32}&C_{33}\end{pmatrix},
 \qquad C_{32}=C_{23}\T.
\end{equation}
For example, $C_{23}$ has one row for each two-point observable and one
column for each three-point observable. It records their correlated
fluctuations. It is not a new physical signal added to the mean.

\paragraph{Subtract the part already predicted by two-point fluctuations.}
Let $\Delta_2$ and $\Delta_3$ be the differences between two chosen means.
The linear prediction of the second difference from the first is
\begin{equation}
 \Delta_{3,\mathrm{pred}}=C_{32}C_{22}^{-1}\Delta_2.
\end{equation}
Subtract it to obtain the remaining difference:
\begin{equation}
 r=\Delta_3-C_{32}C_{22}^{-1}\Delta_2.
\end{equation}
The corresponding covariance is obtained by applying the same subtraction
to random fluctuations. Expanding that covariance gives
\begin{equation}
\begin{split}
 \operatorname{Cov}(d_3-C_{32}C_{22}^{-1}d_2)
 ={}& C_{33}-C_{32}C_{22}^{-1}C_{23}
      -C_{32}C_{22}^{-1}C_{23}\\
    &+C_{32}C_{22}^{-1}C_{22}C_{22}^{-1}C_{23}.
\end{split}
\end{equation}
The last term cancels one of the two subtractions. The remaining matrix is
\begin{equation}
 S=C_{33}-C_{32}C_{22}^{-1}C_{23}.
\end{equation}
It is the conditional covariance of the three-point information after
conditioning on the two-point measurements in a Gaussian model.

\paragraph{A numerical example with only two observations.}
Take the dimensionless covariance and mean difference
\begin{equation}
 \Cov=\begin{pmatrix}4&2\\2&5\end{pmatrix},
 \qquad \Delta=\begin{pmatrix}2\\3\end{pmatrix}.
\end{equation}
The predicted second difference is $(2/4)2=1$. Therefore the residual is
$3-1=2$. Its variance is
\begin{equation}
 S=5-\frac{2\times2}{4}=4.
\end{equation}
The first measurement contributes $2^2/4=1$ to the squared distance; the
remaining second measurement contributes $2^2/4=1$. Thus
\begin{equation}
 \Delta\chi^2_{\rm joint}=1+1=2.
\end{equation}
For comparison, treating the measurements as independent would give
$2^2/4+3^2/5=2.8$, a different answer. The original joint inverse confirms
our subtraction:
\begin{equation}
 \Cov^{-1}=\frac1{16}\begin{pmatrix}5&-2\\-2&4\end{pmatrix}.
\end{equation}
Multiplying this matrix by $\Delta$ gives $(1/4,1/2)^T$. Taking its scalar
product with $\Delta$ gives $2/4+3/2=2$.

\paragraph{The identity tested in the notebook.}
For the complete vector, the same operations give
\begin{equation}
 \Delta\chi^2_{\rm joint}
 =\Delta_2\T C_{22}^{-1}\Delta_2+r\T S^{-1}r.
\end{equation}
The native data covariance stays fixed throughout this comparison. These
squared distances describe fixed means. They are neither marginalized
parameter errors nor detection significances after a fit.

\section{From derivatives to a Fisher forecast}
\paragraph{What one column means.}
Near a fiducial parameter vector $p_0$, approximate the prediction by
\begin{equation}
 \mu(p_0+\delta p)\simeq\mu(p_0)+D\delta p.
\end{equation}
Column $\alpha$ of $D$ is the change of every data entry per unit change of
parameter $p_\alpha$. The notebook evaluates each column with a centered
five-point derivative, using absolute displacements $h_\alpha$:
\begin{equation}
 D_\alpha\simeq
 \frac{\mu(p_0-2h_\alpha e_\alpha)-8\mu(p_0-h_\alpha e_\alpha)
 +8\mu(p_0+h_\alpha e_\alpha)-\mu(p_0+2h_\alpha e_\alpha)}{12h_\alpha}.
\end{equation}
The vector $e_\alpha$ has one in position $\alpha$ and zero elsewhere.
Repeating at half the displacement tests the local differentiation. All
coordinates, including negative $A_1$ and zero $w_a$, retain these absolute
step units in the shared \texttt{differentiate} routine.

\paragraph{Build information before reading errors.}
The likelihood Fisher matrix is
\begin{equation}
 F_{\rm like}=D\T\Cov^{-1}D.
\end{equation}
A Gaussian prior with standard deviation $s_\alpha$ contributes
$1/s_\alpha^2$ to the corresponding diagonal. Hence
\begin{equation}
 F=F_{\rm like}+F_{\rm prior}.
\end{equation}
Both the two-point and joint forecasts receive the same prior once. A flat
prior boundary supplies no local Gaussian curvature. Its possible effect
on a posterior near a boundary is absent from this local calculation.

\paragraph{The new three-point derivative.}
Apply the previous residual subtraction to every derivative column:
\begin{equation}
 E=D_3-C_{32}C_{22}^{-1}D_2.
\end{equation}
The likelihood information then decomposes as
\begin{equation}
 F_{\rm joint,like}=D_2\T C_{22}^{-1}D_2+E\T S^{-1}E.
\end{equation}
The second term is positive semidefinite when $S$ is positive definite.
Thus adding these data cannot reduce exact Fisher information for a fixed
parameter space and covariance. The angle between two derivative columns
can nevertheless become smaller: information growth and a particular
pair's geometric angle are different statements.

\paragraph{How the derivative plots keep their physical meaning.}
The notebook shows a chosen parameter displacement times its derivative,
divided by a row's standard deviation:
\begin{equation}
 t_{i\alpha}=\frac{\Delta p_\alpha D_{i\alpha}}{\sqrt{C_{ii}}}.
\end{equation}
A value of one means one marginal row error. The display displacement is
neither a numerical derivative step nor a prior. Rows are still correlated.
To fit response $b$ to response $a$, minimize
\begin{equation}
 Q(c)=(a-cb)\T\Cov^{-1}(a-cb).
\end{equation}
Differentiating with respect to $c$ gives
\begin{equation}
 \frac{dQ}{dc}=-2b\T\Cov^{-1}a+2c\,b\T\Cov^{-1}b.
\end{equation}
Setting this derivative to zero yields
\begin{equation}
 c=\frac{b\T\Cov^{-1}a}{b\T\Cov^{-1}b}.
\end{equation}
The residual $a-cb$ shows what this one competitor cannot reproduce.
A negative fitted coefficient can indicate an equally strong degeneracy:
opposite-sign responses can be exchanged by opposite parameter changes.
This two-column exercise holds all other parameters fixed.

\section{Marginalization, fixed dark energy and the amplitude combination}
\paragraph{Invert the requested experiment.}
For a resolved Fisher matrix, the parameter covariance is
\begin{equation}
 P=F^{-1}.
\end{equation}
The marginalized error of parameter $\alpha$ is $\sqrt{P_{\alpha\alpha}}$.
The quantity $1/\sqrt{F_{\alpha\alpha}}$ instead fixes every other
parameter. To hold only $w_0,w_a$ fixed, remove their rows and columns from
$F$ before inversion. Deleting them from $P$ would preserve their
marginalization. In the executed notebook both the full and the fixed-dark-energy
experiments resolve; an unconstrained retained parameter would be reported
rather than silently removed, and no pseudoinverse replaces a failed
inverse.

\paragraph{Differentiate the derived amplitude explicitly.}
The commonly used amplitude combination is
\begin{equation}
 S_8=\sigma_8\sqrt{\frac{\Omega_m}{0.3}}.
\end{equation}
Differentiate its square-root factor while holding $\sigma_8$ fixed:
\begin{equation}
 \frac{\partial S_8}{\partial\Omega_m}=\frac{S_8}{2\Omega_m}.
\end{equation}
Next hold $\Omega_m$ fixed and differentiate the linear amplitude factor:
\begin{equation}
 \frac{\partial S_8}{\partial\sigma_8}=\sqrt{\frac{\Omega_m}{0.3}}.
\end{equation}
Writing them as $g_\Omega$ and $g_\sigma$, propagate the covariance:
\begin{equation}
 \operatorname{Var}(S_8)=g_\Omega^2P_{\Omega\Omega}
 +2g_\Omega g_\sigma P_{\Omega\sigma}+g_\sigma^2P_{\sigma\sigma}.
\end{equation}
For a teaching example take $\Omega_m=0.3$, $\sigma_8=0.8$, and
\begin{equation}
 P_{\Omega,\sigma}=\begin{pmatrix}.0009&-.0024\\-.0024&.01\end{pmatrix}.
\end{equation}
Then $g_\Omega=4/3$ and $g_\sigma=1$. Substituting each term gives
\begin{equation}
 \operatorname{Var}(S_8)=.0016-.0064+.01=.0052.
\end{equation}
Dropping the covariance would give $.0116$, more than twice the variance.
These numbers are an algebra test, not a DES result.

\paragraph{Read a figure of merit only after marginalizing.}
Transform the full parameter covariance to the chosen two-coordinate
covariance $P_{\rm pair}$. Its figure of merit is
\begin{equation}
 \mathrm{FoM}=\frac1{\sqrt{\det P_{\rm pair}}}.
\end{equation}
Compare the same pair, parameter space, fiducial, priors and cuts between
experiments. The notebook's NLA and TATT forecasts vary different IA
spaces; their cosmological comparison includes that change in
marginalization. Source-calibration prior widths are listed explicitly in
the notebook, and there are no added Gaussian priors on $w_0,w_a$.

\section{What must be checked before interpreting a result}
\paragraph{Derivatives need more than a stable mean.}
A small error in two nearby predictions can become large after division by
a small displacement. Compare both the half-step derivatives and complete
Fisher matrices at finer numerical sampling. The associated accuracy
notebook also uses the shared fourth-difference noise screen. Its
independent-noise assumptions do not detect every coherent derivative
error. A stable mean, a local noise estimate and a converged marginalized
FoM answer different numerical questions.

\paragraph{Keep the published-analysis differences visible.}
The numbered examples currently use a fixed Gaussian interpretation of a
simulation covariance and their declared source preparation. These Fisher
experiments neither reproduce a finite-simulation posterior correction nor
automatically reproduce the paper's compression and scale selection.
Changing the data mask requires selecting the covariance and row metadata
in the same order.

\paragraph{The executed forecasts.}
For the finite comparison of Section 2, the opposing NLA means differ by
$\Delta\chi^2=15.06$ jointly, splitting into 7.83 from the two-point rows
and 7.24 from the conditional three-point residual; the opposing TATT
means give 7.48, splitting into 5.12 and 2.36. Halving every derivative
step changes the NLA derivative columns by at most 0.074\%. The signed
response cosines between the alignment-amplitude and sampled-amplitude
columns are $-0.75$ (two-point) and $-0.80$ (joint) for NLA, and 0.74 and
0.80 for TATT. Adding the aperture rows tightens the alignment coordinates: the
joint-to-two-point marginalized error ratios are 0.90 for $A_1$ and 0.61
for $\eta_1$ in NLA, and 0.88, 0.54, 0.64, 0.83 and 0.81 for
$A_1,\eta_1,A_2,\eta_2,b_{\rm TA}$ in TATT. With free dark energy, $\sigma(S_8)$ improves from
0.115 to 0.060 (NLA) and from 0.128 to 0.081 (TATT), and the
$(\Omega_m,S_8)$ figure of merit from 150 to 458 and from 100 to 340.
With $w_0,w_a$ fixed, the figures of merit become 915 to 1914 (NLA) and
636 to 1805 (TATT). These are local Fisher statements under this note's
fixed covariance and parameter spaces, not reproductions of either
paper's posterior.

\paragraph{Acknowledgements.}
\textbf{We thank Sunao Sugiyama and Rafael Gomes for their codes and scientific
work, which made these comparisons possible.} The upstream code references are
\href{https://github.com/git-sunao/fastnc}{fastnc},
\href{https://github.com/rchgomes/3pcf_integrator}{3pcf\us integrator} and
\href{https://github.com/rchgomes/desy6-3pcf}{desy6-3pcf}; this project
compares with saved fixtures from those codes and does not run them.
CosmoLike supplies the projection and Fisher machinery; the companion
notebook follows Cocoa's ordinary Cobaya likelihood path.

\begin{thebibliography}{9}
\bibitem{methods} R. C. H. Gomes et al., \emph{Cosmology with second and third-order
shear statistics for the Dark Energy Survey: Methods and simulated analysis},
\href{https://arxiv.org/abs/2503.03964}{arXiv:2503.03964}.
\bibitem{tatt} R. C. H. Gomes et al., \emph{Tidal alignment and tidal torquing modeling
for the cosmic shear three-point correlation function and mass aperture skewness},
\href{https://arxiv.org/abs/2601.09133}{arXiv:2601.09133}.
\end{thebibliography}
\end{document}
